\section{Results}

\subsection{Evidence Base}

The systematic review retained 17 sources for evidence extraction and thematic synthesis following application of inclusion and exclusion criteria. These included 6 systematic or umbrella reviews, 4 scoping or narrative reviews, 4 government and institutional reports (including GAO audit findings and ONC reports to Congress), and 3 foundational framework and policy analyses. Additionally, 4 national eHealth implementation cases were analyzed using primary source documentation. The full evidence extraction matrix, including source-level U-T-I-O domain coding and key findings, is provided in Table~\ref{tab:evidence}.

\subsection{Systematic Review Findings}
Across all 17 retained sources, adoption determinants consistently cluster into four domains corresponding to the U-T-I-O Model. No single technical or organizational variable emerged as independently determinative of adoption outcomes; rather, determinants across domains consistently co-occur and interact in shaping implementation success or failure.

User readiness determinants include usability, workflow fit, trust, training, perceived usefulness, and user satisfaction. Technical determinants include information quality, system quality, semantic and syntactic interoperability, reliability, interface integration, and data completeness. Institutional determinants include regulatory incentives, national digital health strategy, privacy law, interoperability mandates, and information governance. Organizational determinants include leadership, implementation planning, stakeholder engagement, change management, vendor governance, and resource capacity.

\subsection{Multi-Case Findings}
Four public cases were selected to compare determinant patterns across different implementation environments: ONC-supported U.S. health information exchange programs, NHS Shared Care Records, Estonia's national eHealth system, and VA/DoD federal EHR modernization.

\begin{table}[H]
\centering
\begin{tabular}{lccccp{5.2cm}}
\toprule
Case & U & T & I & O & Interpretation \\
\midrule
ONC HIE & Medium & High & High & Medium & Strong policy and interoperability infrastructure; variable organizational uptake. \\
NHS Shared Care Records & Medium & High & High & Medium & Strong national interoperability vision; local implementation and information governance remain central. \\
Estonia eHealth & High & High & High & High & Mature national architecture with legal mandate and broad provider participation. \\
VA/DoD EHR & Medium & Medium & High & Low--Medium & Institutional mandate is strong, but configuration, usability, and execution risks persist. \\
\bottomrule
\end{tabular}
\caption{Comparative Case Analysis of Adoption Determinants}
\label{tab:case-comparison}
\end{table}

\subsection{Synthesis}

The systematic review and multi-case findings converge on a consistent pattern. The evidence supports the proposed model:

\begin{equation}
Adoption = f(U, T, I, O)
\end{equation}

Across the literature, the four U-T-I-O domains consistently appear as co-determinants of adoption outcomes. No reviewed source identified a single-domain explanation as sufficient. Cases in which one or more domains were weak—regardless of the strength of others—exhibited measurably lower adoption value, higher user resistance, or implementation delays.

The case comparison reinforces this finding empirically. Technical interoperability is necessary but demonstrably insufficient: both ONC HIE programs and the VA/DoD initiative deployed substantial technical infrastructure, yet neither achieved the breadth of adoption or user engagement observed in Estonia's nationally aligned system. The decisive differentiator across cases is multi-domain convergence. Estonia's high ratings across all four domains correspond to the highest assessed adoption maturity. The VA/DoD case, despite strong institutional mandate (I) and significant technical investment (T), exhibits the lowest overall adoption value due to persistent gaps in user readiness (U) and organizational execution (O)—an outcome directly consistent with the U-T-I-O Model's prediction that domain-level misalignment constrains adoption regardless of isolated domain strengths.

These findings validate the U-T-I-O Model as both a descriptive and diagnostic tool: it accurately characterizes observed adoption patterns and provides an explanatory account of why well-funded, mandate-backed implementations can still fall short of their intended outcomes.
